一摩尔流体满足范德瓦耳斯方程 \(P=RT/(v-b)-a/v^{2}\)，a,b,R>0，v>b。假设临界点是同一等温曲线上同时满足 \(\partial P/\partial v=0\)、\(\partial ^{2}P/\partial v^{2}=0\) 的点。
(1) 联立临界条件求临界摩尔体积、温度和压强；(2) 求临界压缩因子 Zc=Pcvc/(RTc)，与理想气体比较；(3) 从等温压缩系数 \(\kappa T=-(1/v)(\partial v/\partial P)T\) 讨论机械稳定区间；(4) 沿 Tc 等温线作 \(v-vc\) 小偏差展开至首次非零阶，说明临界处为何压力变化是三次而非一次。
